% Quantifying Attention Flow in Transformers (ACL 2020)
% Authors: Samira Abnar, Willem Zuidema (University of Amsterdam)
% Consolidated technical content

\section{Abstract}
In the Transformer model, "self-attention" combines information from attended embeddings into the representation of the focal embedding in the next layer. Across layers, information from different tokens gets increasingly mixed, making attention weights unreliable as explanation probes. We propose two methods for approximating attention to input tokens: Attention Rollout and Attention Flow, as post hoc methods when using attention weights as relative relevance of input tokens. Both methods give complementary views on information flow and yield higher correlations with importance scores from ablation and input gradients compared to raw attention.

\section{Problem Statement}
As we go deeper into the model, embeddings become more contextualized and may all carry similar information. Raw attention weights in higher layers become uniform and lack interpretable patterns. This underscores the need to track attention weights back to the input layer.

\section{Methodology}

\subsection{Modeling Information Flow}
Model information flow in the network with a DAG (Directed Acyclic Graph):
- Nodes: input tokens and hidden embeddings
- Edges: attentions from nodes in each layer to previous layer
- Weights: attention weights

\subsection{Accounting for Residual Connections}
Residual connections play significant role in tying corresponding positions in different layers. Given attention module with residual connection:
\begin{equation}
V_{l+1} = V_l + W_{att}V_l = (W_{att} + I)V_l
\end{equation}

To account for residual connections, add identity matrix and re-normalize:
\begin{equation}
A = 0.5W_{att} + 0.5I
\end{equation}

\subsection{Attention Rollout}
Track information propagation from input tokens to higher layer embeddings. If we view edge weights as proportion of information transferred, compute total information propagated by summing over all possible paths.

Recursive computation:
\begin{equation}
\tilde{A}(l_i) = \begin{cases}
A(l_i)\tilde{A}(l_{i-1}) & \text{if } i > j \\
A(l_i) & \text{if } i = j
\end{cases}
\end{equation}

where $\tilde{A}$ is Attention Rollout, $A$ is raw attention, and multiplication is matrix multiplication.

\subsection{Attention Flow}
Treat attention graph as a flow network:
- Capacity: attention weights on edges
- Source: hidden embeddings
- Sink: input tokens

Use maximum flow algorithm to compute maximum attention flow from any node to input nodes.

\textbf{Key difference from Rollout:}
- Weight of single path = minimum value of edge weights (not product)
- Cannot simply sum path weights due to potential overlap/overflow

\subsection{Complexity}
- Attention Rollout: $O(d \cdot n^2)$
- Attention Flow: $O(d^2 \cdot n^4)$
where $d$ = model depth, $n$ = number of tokens

\section{Analysis}

\subsection{Visualization Findings}
- First layer of Rollout and Flow are same (only differ by residual connections from raw)
- As we move to higher layers, residual connections fade away
- Unlike raw attention, Rollout and Flow patterns become MORE distinctive in higher layers

\subsection{Key Differences}
\textbf{Attention Flow:}
- Weights amortized among most attended tokens
- Indicates SET of important input tokens
- More relaxed assumptions lead to more accurate results

\textbf{Attention Rollout:}
- Weights more focused
- Stricter assumptions
- Views weights as proportion factors

\section{Experiments}

\subsection{Task}
Verb number prediction (singular/plural) on subject-verb agreement dataset. Model: 6-layer, 8-head Transformer with 128 hidden size. Accuracy: 0.96

\subsection{Evaluation Methods}
\textbf{Blank-out:} Replace each token with UNK, measure effect on correct class probability
\textbf{Input Gradients:} Gradient-based importance scores

\subsection{Quantitative Results}
SpearmanR correlation with blank-out scores (2000 samples):

| Method | L1 | L2 | L3 | L4 | L5 | L6 |
|--------|-----|-----|-----|-----|-----|-----|
| Raw | 0.69 | 0.10 | -0.11 | -0.09 | 0.20 | 0.29 |
| Rollout | 0.32 | 0.38 | 0.51 | 0.62 | 0.70 | 0.71 |
| Flow | 0.32 | 0.44 | 0.70 | 0.70 | 0.71 | 0.70 |

Key finding: Both Rollout and Flow have INCREASING correlation with importance as layer depth increases, while Raw attention correlation degrades in middle layers.

\subsection{BERT Experiments}
- DistillBERT on SST-2: All methods have low correlation with gradients, but Rollout/Flow still better than Raw
- Pronoun resolution: Flow weights most consistent with model predictions

\section{Conclusions}
1. Translating embedding attentions to token attentions provides better explanations
2. Methods are simple, task/architecture agnostic, require only attention weights
3. For Transformer decoders (with causal masking), normalize based on receptive field first
4. Future work: use effective attention weights or gradient-based adjustments
